\chapter{Introduction} \label{cap:introduction}

This report describes the testing work performed for the Calendar Spring Boot application supplied for VVS Assignment 2. The assignment goal is to assess the ability to test a larger Spring Boot MVC application at several levels: business logic, third-party integrations, REST/API behavior, persistence, and browser-level workflows.

The application is a meeting scheduler. Users can register and log in, create meetings, invite other users, accept or decline invitations, expose an iCalendar feed, and copy discovered events from external providers into their own calendar.

The project repository is available at \url{https://github.com/DanielSousa01/calendar}.

This document is written as a walkthrough of the solution. Chapter~\ref{cap:sut} first explains the system under test and the modification policy. Chapter~\ref{cap:strategy} maps each assignment requirement to the implemented evidence. The following chapters then walk through each testing level, explaining what part of the code is exercised, why that testing level was chosen, what data or mocks are used, and which defects are intentionally exposed.

The test suite was designed with three goals:

\begin{itemize}
\item Cover each testing category requested in the assignment using the most appropriate isolation level.
\item Keep the application implementation unchanged, so the tests evaluate the code as it exists in the submitted project.
\item Expose implementation bugs through executable tests, while keeping the default CI suite deterministic and passing.
\end{itemize}

The default Maven run executes 38 non-bug-revealing tests and passes. The remaining nine bug-revealing tests are tagged with \texttt{@Tag("bug")} and can be run deliberately with \texttt{mvn -Pbug-tests test}, they fail because they reveal current implementation defects documented in Chapter~\ref{cap:bugs}.
